\subsection{Valor de nuestro trabajo}

%\textbf{NOTA}: Aquí se mencionan los aportes que puede hacer nuestro trabajo y cual es su valor científico, social y cultural.

Este trabajo busca plantar un antecedente sobre los proyectos que se han realizado alrededor del aprovechamiento de los desechos orgánicos, integrando y enriqueciendo diversas ideas que se han realizado por separado, como se vieron en la tabla de antecedentes (Tabla \ref{tab:Tab_Antecedentes}).

De esta manera, se obtiene un sistema mas robusto, especializado para el tratamiento de los residuos vegetales y frutales, cuidando las condiciones del ambiente donde se realiza la maduración de los residuos, obteniendo resultados en 72 horas, que en otros tipos de procesos de compostaje, se obtienen en 2 meses.

Las dimensiones de este prototipo se pueden escalar, para aumentar su capacidad y poder beneficiar a diversas poblaciones, como lo pueden zonas habitacionales, edificios de departamentos, escuelas e incluso mercados. 

Aprovechando este tipo de recursos, se beneficia al ambiente, evitando la contaminación producida por el transporte y la descomposición de los desechos en  sanitarios. De igual forma, se busca cerrar el ciclo natural, devolviendo los nutrientes generados por las plantas y arboles, al suelo, para seguir produciendo vida.

